% ============================================================
% PAGE 7 — Module 3: Detection Head and Optimized Loss Functions
% ============================================================

\subsection{Module 3: Detection Head and Optimized Loss
Functions}

The detection head is the terminal stage of the pipeline---the
point at which abstract, multi-scale feature representations
are decoded into the concrete set of bounding boxes, class
labels, and confidence scores that constitute the detector's
output.  FusionDet employs a \textit{decoupled head} topology
in which classification and localization are handled by
separate, parallel sub-networks rather than a single shared
branch.  This decoupling, first systematized by
YOLOX~\cite{ge2021yolox}, resolves a long-standing conflict in
single-stage detector design: classification benefits from
features that emphasize semantic discriminability, while
regression benefits from features that preserve precise spatial
gradients.  Forcing both tasks through a shared representation
produces a compromise that is optimal for neither.

\subsubsection{Head Architecture}

For each pyramid level $l \in \{2,3,4,5\}$, the recalibrated
feature map $\hat{\mathbf{N}}_{l}$ produced by the neck is fed
into two lightweight branches operating in parallel:

\textit{Classification branch.}  Two stacked $3\times3$
depthwise-separable convolutions, each followed by batch
normalization and SiLU activation, reduce the feature map to a
tensor $\mathbf{O}_{\text{cls}}^{l} \in
\mathbb{R}^{K \times H_l \times W_l}$, where $K$ is the number
of object categories.  A sigmoid activation is applied
element-wise to produce per-class probability maps.

\textit{Regression branch.}  An identically structured pair of
depthwise-separable convolutions produces a tensor
$\mathbf{O}_{\text{reg}}^{l} \in
\mathbb{R}^{4 \times H_l \times W_l}$ encoding the four
bounding-box parameters $(x_c, y_c, w, h)$ in the center-size
format defined in (\ref{eq:bbox}).

\textit{Objectness branch.}  A single $1\times1$ convolution
followed by sigmoid activation outputs a scalar objectness
score $\mathbf{O}_{\text{obj}}^{l} \in
\mathbb{R}^{1 \times H_l \times W_l}$ at each spatial
location, indicating the likelihood that the location
corresponds to the center of any object regardless of class.

At inference time the final detection confidence for class $k$
at location $(i,j)$ of level $l$ is the product:
\begin{equation}
    \hat{s}_{k,i,j}^{\,l} =
    \mathbf{O}_{\text{obj}}^{l}(i,j) \;\cdot\;
    \mathbf{O}_{\text{cls}}^{l}(k,i,j)
    \label{eq:conf}
\end{equation}
Predictions from all four levels are concatenated, ranked by
$\hat{s}$, and filtered through Non-Maximum Suppression (NMS)
with an IoU threshold of $0.65$ to produce the final detection
set $\mathcal{D}$.

\subsubsection{Multi-Part Loss Formulation}

Training the three-branch head requires a composite loss that
simultaneously supervises classification accuracy, bounding-box
quality, and objectness calibration.  The total training
objective is:
\begin{equation}
    \mathcal{L}_{\text{total}} =
    \lambda_{\text{cls}}\,\mathcal{L}_{\text{cls}}
    \;+\;
    \lambda_{\text{box}}\,\mathcal{L}_{\text{CIoU}}
    \;+\;
    \lambda_{\text{obj}}\,\mathcal{L}_{\text{obj}}
    \label{eq:total_loss}
\end{equation}
where $\lambda_{\text{cls}}$, $\lambda_{\text{box}}$, and
$\lambda_{\text{obj}}$ are scalar coefficients controlling the
relative contribution of each term.  We set
$\lambda_{\text{cls}} = 1.0$, $\lambda_{\text{box}} = 5.0$, and
$\lambda_{\text{obj}} = 1.0$ following the weighting convention
of~\cite{bochkovskiy2020yolov4}, validated through a
hyperparameter sweep in our ablation study.

\textbf{Classification loss.}  We adopt the binary
cross-entropy variant of focal loss~\cite{lin2017focal} to
address the severe foreground--background imbalance:
\begin{equation}
    \mathcal{L}_{\text{cls}} =
    -\frac{1}{N_{\text{pos}}}
    \sum_{i}\;
    \alpha_{t}\,(1 - p_{t,i})^{\gamma}\,
    \log(p_{t,i})
    \label{eq:focal}
\end{equation}
where $p_{t,i}$ is the predicted probability for the true class
of sample $i$, $\alpha_{t}$ is a per-class balancing factor,
and $\gamma$ is the focusing parameter.  Setting $\gamma = 2.0$
and $\alpha_{t} = 0.25$ causes the loss to down-weight the
overwhelming majority of easy negative locations, concentrating
gradient energy on the sparse, hard positives.

\textbf{Objectness loss.}  The objectness branch is supervised
with binary cross-entropy:
\begin{equation}
    \mathcal{L}_{\text{obj}} =
    -\frac{1}{N}
    \sum_{i}\Bigl[
        y_{i}\log(\hat{o}_{i})
        + (1-y_{i})\log(1-\hat{o}_{i})
    \Bigr]
    \label{eq:obj}
\end{equation}
where $\hat{o}_{i}$ is the predicted objectness score and
$y_{i} \in \{0,1\}$ is the ground-truth label assigned via
the SimOTA dynamic label assignment strategy~\cite{ge2021yolox}.

\textbf{Bounding-box regression loss --- CIoU.}  FusionDet
employs the \textit{Complete IoU} (CIoU) loss~\cite{zheng2020distance},
which extends the basic IoU formulation along three
complementary axes: overlap area, center-point distance, and
aspect-ratio consistency:
\begin{equation}
    \mathcal{L}_{\text{CIoU}} =
    1 - \text{IoU}
    + \frac{\rho^{2}(\mathbf{b}_{p},\,\mathbf{b}_{gt})}{c^{2}}
    + \alpha_{\text{ar}}\,v
    \label{eq:ciou}
\end{equation}

\textit{Distance term.}
$\rho(\mathbf{b}_{p},\mathbf{b}_{gt})$ is the Euclidean
distance between the predicted and ground-truth box centers,
and $c$ is the diagonal length of the smallest enclosing
rectangle:
\begin{equation}
    \frac{\rho^{2}(\mathbf{b}_{p},\,\mathbf{b}_{gt})}{c^{2}}
    = \frac{(x_c^{p} - x_c^{gt})^{2} +
            (y_c^{p} - y_c^{gt})^{2}}
           {c_{w}^{2} + c_{h}^{2}}
    \label{eq:dist}
\end{equation}
This normalized penalty accelerates convergence when two boxes
share little or no overlap, a regime in which the gradient of
the raw IoU loss vanishes.

\textit{Aspect-ratio consistency.}  The parameter $v$
quantifies aspect-ratio discrepancy:
\begin{equation}
    v = \frac{4}{\pi^{2}}
    \left(
        \arctan\frac{w^{gt}}{h^{gt}}
        - \arctan\frac{w^{p}}{h^{p}}
    \right)^{\!2}
    \label{eq:v}
\end{equation}
with the adaptive trade-off coefficient:
\begin{equation}
    \alpha_{\text{ar}} = \frac{v}{(1 - \text{IoU}) + v}
    \label{eq:alpha_ar}
\end{equation}
so that aspect-ratio correction receives increasing weight as
the overlap improves---a principled schedule that prevents the
$v$ term from dominating early in training when the boxes are
still poorly localized.

Together, the three CIoU components ensure informative,
non-vanishing gradients across the full spectrum of prediction
quality, from grossly misaligned boxes to nearly perfect
localizations, yielding faster convergence and tighter final
bounding boxes than either smooth-$L_1$ or basic IoU loss
alone~\cite{zheng2020distance}.
